\bookchapter{第二十七章：爸媽也能輕鬆用——適老化手機完全指南}

\chapterintro{說句實話，寫這一章的時候，我想到了我爸媽。我媽每次手機彈出一個“系統更新”，她都不敢點，怕把手機弄壞了。我爸的微信字型比螞蟻還小，但他不知道怎麼調大，就這麼眯著眼睛看了兩年。這不是他們“笨”，是手機預設把所有人都當成了二十多歲的年輕人。這一章，我用最直白的話，把長輩們手機使用中最頭疼的問題一個個解開。}

\booksection{零、先把心態擺正：這不是“老人用不來”，是“手機沒做好”}

在做任何設定之前，有一個觀念要先說清楚：\textbf{現在的智慧手機介面，從設計之初就沒怎麼考慮老年使用者}。彈窗小得看不清、操作手勢需要精準觸碰、圖示長得都差不多——這些不是使用者的問題，是產品的問題。所以帶長輩用手機，第一件事不是教，是幫他們\textbf{把環境調好}。環境調好了，大半的問題就消失了。

帶長輩用手機有三條鐵律：(1)\textbf{別催}——你用了十年的智慧機，他們可能才用了三個月。(2)\textbf{用類比}——“Wi-Fi就像家裡的水管，開了水龍頭才有水”，這種比喻比專業術語好理解一萬倍。(3)\textbf{寫下來}——口頭教一遍他們可能忘了，你拿張紙把關鍵步驟寫下來貼在他手機殼背面，比什麼都管用。

\booksection{一、第一步：開啟“簡易模式”——手機變清爽了}

現在的手機基本都內建了一個“簡易模式”（有的叫“關懷模式”“老年模式”“長輩模式”），開啟之後，整個手機都會大變樣：

\begin{itemize}
  \item \textbf{字型和圖示變得很大}——再也不用戴老花鏡看手機了。
  \item \textbf{桌面圖示大幅精簡}——只保留電話、簡訊、相機、相簿、微信等最常用的幾個App，其他全收起來。
  \item \textbf{操作變簡單了}——很多花裡胡哨的功能（負一屏、智慧助手、浮窗等）預設關閉，介面主打一個清爽。
  \item \textbf{聲音變大了}——來電鈴聲、通知音量、通話聲音都自動調到最大。
\end{itemize}

\booksubsection{各品牌簡易模式開啟方法（直接照著操作就行）：}

\begin{center}
\begin{tabularx}{\linewidth}{|X|X|X|}
\hline
\textbf{品牌} & \textbf{模式名稱} & \textbf{開啟方法} \tabularnewline
\hline
iPhone & 輔助使用（沒有單獨簡易模式，但可逐項調大） & 設定 > 顯示與亮度 > 文字大小，拉到最大。再進"輔助功能 > 顯示與文字大小"開"粗體文字""按鈕形狀" \tabularnewline
華為/榮耀 & 簡易模式 & 設定 > 系統和更新 > 簡易模式 > 開啟 \tabularnewline
小米/紅米 & 極簡模式 & 設定 > 更多設定 > 極簡模式 > 開啟 \tabularnewline
OPPO/一加 & 簡易模式 & 設定 > 特色功能 > 簡易模式 \tabularnewline
vivo/iQOO & 簡易桌面 & 設定 > 桌面、鎖屏與桌布 > 簡易桌面 \tabularnewline
三星 & 簡易模式 & 設定 > 顯示 > 簡易模式 \tabularnewline
\hline
\end{tabularx}
\end{center}

如果你用的是其他品牌，在手機“設定”頂部的搜尋框裡直接搜“簡易”或“長輩”兩個字，一般都能找到。找不到的話，按下面的“手動調大”方案一條條設定，效果也一樣。

\booksection{二、如果不開簡易模式，手動調也可以——“大字+大圖示”五步走}

有些長輩不想用簡易模式，覺得介面太“老氣”。那就手動調以下五項，效果一樣好：

\begin{enumerate}
  \item \textbf{字型調到最大}：設定 > 顯示 > 字型大小，拖到最右邊。順便把"顯示大小"（有的叫"顯示縮放"）也調到最大，這樣所有介面元素都跟著變大。
  \item \textbf{圖示變大}：在主螢幕雙指捏合（像縮小照片那樣）進入桌面設定 > 桌面佈局（或圖示大小），選格子最少的佈局（比如3×5），圖示自然就大了。
  \item \textbf{開啟"高對比度文字"}：讓文字邊緣更銳利，在"輔助功能"裡可以找到這個開關。
  \item \textbf{關閉不必要的通知}：設定 > 通知，把除了電話、簡訊、微信之外的App通知全關掉。這樣手機不會一會兒響一下，不會讓人覺得煩。
  \item \textbf{把常用聯絡人放到桌面}：長按桌面空白處 > 新增小工具 > 聯絡人，把老伴、兒女、醫生、緊急聯絡人的快捷撥號放到桌面上。點一下就撥出去，不用翻通訊錄。
\end{enumerate}

\begin{tcbpractice}{動手：幫長輩"裝修"一下手機}
\begin{enumerate}
  \item 開啟簡易模式（按上面的品牌表格操作）。
  \item 把電量數字顯示開啟（設定 > 電池 > 電量百分比），讓他們隨時知道還剩多少電。
  \item 調出通知欄，把藍芽、NFC、飛航模式這些長輩用不著的快捷開關移除或藏起來。
  \item 在桌面上放四個快捷撥號：老伴、兒女、家庭醫生（如果有）、社群/物業。
\end{enumerate}

\end{tcbpractice}

\booksection{三、防詐騙——這章請一定帶著長輩一起看}

老年人是電信詐騙的重災區。不是他們警惕性低，而是騙子的話術經過了精心設計，專門針對老年人的心理弱點。以下是我整理的最常見騙局型別和應對方法，建議你直接念給長輩聽，或者列印出來貼在電話旁邊。

\booksubsection{六大常見騙局（看到了立刻掛電話）：}

\begin{itemize}
  \item \textbf{(1) "你孫子出事了，趕緊匯錢"}——騙子會說"你孫子/孫女在外面闖禍了，被警察抓了，需要保證金"。記住：\textbf{任何要求轉帳才能"放人"的電話都是詐騙}。公檢法不會電話要錢。
  \item \textbf{(2) "你的醫保卡被停用了"}——簡訊裡說醫保卡異常，讓你點連結"啟用"。記住：\textbf{醫保局、社保局不會發簡訊讓你點連結}。遇到這種情況直接打12333（社保官方熱線）核實。
  \item \textbf{(3) "高收益理財，穩賺不賠"}——微信群、簡訊裡推薦各種"養老理財""國家扶持專案"，承諾高額回報。記住：\textbf{承諾年化收益超過5\%還"保本"的，99.9\%是騙局}。
  \item \textbf{(4) "我是你領導/老朋友，換號碼了"}——冒充熟人借錢。記住：\textbf{打電話核實比什麼都管用}，說"不方便接電話"的八成是騙子。
  \item \textbf{(5) "免費領雞蛋/保健品講座"}——線下免費活動誘導老人聽"健康講座"，然後賣天價保健品。記住：\textbf{天上不掉餡餅，免費活動背後都是生意}。
  \item \textbf{(6) "你的手機中毒了，需要遠端修復"}——騙子讓你下載一個遠端控制App（如向日葵、TeamViewer），然後遠端操控你的手機轉走錢。記住：\textbf{絕對不要讓任何人遠端操控你的手機！}手機廠商、通訊運營商、銀行都不會給你打電話讓你裝遠端控制軟體。
\end{itemize}

\booksubsection{三道防線，立好之後能擋掉90\%的詐騙：}

\begin{enumerate}
  \item \textbf{第一道：手機自帶攔截}——安卓手機開啟"騷擾攔截"（在撥號App或手機管家App裡），iPhone開啟"靜音未知來電"（設定 > 電話 > 靜音未知來電）。
  \item \textbf{第二道：運營商攔截}——傳送簡訊"KTFSR"到10086，免費開通移動的高頻騷擾電話攔截；電信聯通也有類似服務。打運營商客服電話讓人工幫你開通就行。
  \item \textbf{第三道：國家反詐中心App}——在應用商店下載"國家反詐中心"App，註冊並開啟來電預警。這個App和公安的資料庫聯動，騙子號碼一撥過來就能識別。
\end{enumerate}

\begin{tcbwarn}{給長輩的"防騙口訣"（列印出來貼電話旁邊）}
\begin{itemize}
  \item \textbf{陌生電話要轉帳 掛掉！}
  \item \textbf{簡訊連結要你點 刪掉！}
  \item \textbf{說高收益保本的 拉黑！}
  \item \textbf{讓你裝遠端控制軟體 直接掛！}
  \item \textbf{真警察不會電話辦案更不會要錢！}
  \item \textbf{拿不準的事，先給兒女打個電話問問。}
\end{itemize}

\end{tcbwarn}

\booksection{四、遠端協助——你在外地也能幫爸媽修手機}

有時候長輩遇到的問題，你在電話裡描述半天他們也找不到對應的按鈕。這時候“遠端協助”功能就派上用場了——你可以直接在自己的手機上看到他們手機的螢幕，幫他們操作。各品牌的方案不一樣：

\begin{itemize}
  \item \textbf{華為/榮耀}：自帶"暢連"App，裡面有"螢幕共享"功能。你和爸媽都裝暢連，發起視訊通話後點選"共享螢幕"，你就可以在他們螢幕上畫圈標註，告訴你媽媽"點這裡"。
  \item \textbf{小米/紅米}：小米視訊通話App裡有"遠端協助"，支援螢幕共享+語音通話同步。
  \item \textbf{OPPO/vivo}：內建"家庭空間"或"遠端守護"，在設定裡搜尋"遠端"就能找到。
  \item \textbf{通用方案（所有品牌都可用）}：下載"向日葵遠端控制"App（免費版夠用），你們各裝一個。註冊同一個帳號或加好友，讓爸媽把手機上顯示的驗證碼念給你，你就能遠端看到和操作他們的螢幕了。
  \item \textbf{iPhone使用者}：FaceTime通話中支援"共享螢幕"。在FaceTime通話時，點選螢幕共享按鈕，對方就能看到你的螢幕。但iOS安全限制無法遠端操作，只能"指揮式教學"。
\end{itemize}

\booksection{五、微信的“適老改造”——跟長輩說這幾句就夠了}

微信是長輩們用得最多的App，但介面太密太雜。以下是幾個立竿見影的調整：

\begin{itemize}
  \item \textbf{調大微信自己的字型}：微信 > 我 > 設定 > 通用 > 字型大小，拖到最大。注意，這個隻影響微信內部的字型，不影響手機系統字型。
  \item \textbf{開啟"關懷模式"}：微信 > 我 > 設定 > 關懷模式 > 開啟。開啟後字型和按鈕都會變大，色彩對比度加強，還能一鍵"聽文字訊息"（點一下文字訊息就能朗讀出來）。
  \item \textbf{幫他們備註好重要聯絡人}：把兒女、親戚的名字在通訊錄裡改成容易認的——比如"大兒子-小明""二女兒-小紅"，讓他們一開啟微信就知道誰是誰。
  \item \textbf{教他們用"語音輸入"而不是打字}：聊天框旁邊的鍵盤上，按住空格鍵（安卓）或點麥克風圖示（iPhone），直接說話就能變成文字傳送。識別準確率已經非常高了，對不會拼音打字的老人來說是"質變級"的便利。
  \item \textbf{關閉"發現頁"裡不用的功能}：微信 > 我 > 設定 > 通用 > 發現頁管理，把"朋友圈""影片號""直播""購物""遊戲"裡不用的全關掉，減少幹擾。
\end{itemize}

\booksection{六、手機丟了怎麼辦？幫長輩做個“應急預案”}

預則立，不預則廢。趁手機還沒丟，先把這幾件事做好：

\begin{enumerate}
  \item \textbf{確認"查詢我的手機"已開啟}：iPhone在設定 > 你的名字 > 查詢 > 查詢我的iPhone；安卓在設定搜"查詢裝置"。務必確認這個開關是亮著的。
  \item \textbf{把重要的密碼和帳號記在一張紙上}：Apple ID / 華為帳號 / 小米帳號的密碼，鎖在家裡抽屜裡一份，給信得過的兒女一份。丟了手機之後要登入這些帳號來遠端鎖定和定位，如果忘了密碼就麻煩了。
  \item \textbf{教長輩一句話}：\textbf{"手機丟了，第一時間給家人打電話。"}家人可以立刻登入iCloud或安卓的查詢裝置網頁，遠端標記為"丟失模式"，這樣撿到的人也打不開、用不了。
\end{enumerate}

\booksection{七、無線耳機和助聽器相容——被忽視的好功能}

很多長輩有聽力問題，但不知道手機可以跟助聽器協同工作。iPhone和部分安卓手機支援MFi（Made for iPhone）助聽器，可以直接把電話聲音傳到助聽器裡，比舉著手機聽筒清楚太多了。

\begin{itemize}
  \item \textbf{連線方法}：iPhone設定 > 輔助功能 > 助聽裝置，開啟藍芽，按照助聽器的配對說明操作。
  \item \textbf{如果不用助聽器}：在"輔助功能 > 音訊/視覺"裡開啟"電話降噪"和"單聲道音訊"，能改善通話清晰度。
  \item \textbf{安卓}：在"輔助功能 > 聽力增強"裡有類似設定，各品牌位置略有差異，直接搜"助聽"兩個字。
\end{itemize}

\booksection{八、健康相關功能——手機不只是通訊工具}

\begin{itemize}
  \item \textbf{SOS緊急求助}：連續按電源鍵五次（大多數手機），手機會自動撥打緊急電話並給預設的緊急聯絡人傳送定位。幫長輩提前設定好——iPhone："設定 > 緊急SOS"，安卓："設定 > 安全 > 緊急求助"。教長輩用之前先演習一遍，免得到時候緊張按不準。
  \item \textbf{跌倒檢測}：Apple Watch有這個功能，檢測到你摔倒且一段時間不動，會自動撥打急救電話並通知家人。有預算的話給長輩配一個，關鍵時刻能救命。
  \item \textbf{用藥提醒}：iPhone的"健康"App內建用藥提醒功能，設好藥名、時間和劑量，到點就彈出提醒。安卓可以用"定時鬧鐘"或下載"吃藥提醒"類App。
\end{itemize}

\begin{tcbsummary}{本章小結}
\begin{itemize}
  \item 先開簡易模式或手動調大字型/圖示，環境調好大半問題就沒了
  \item 帶長輩用手機三鐵律：別催、用類比、寫下來
  \item 六大常見騙局要反覆提醒，三道防線立好能擋掉90\%詐騙
  \item 遠端協助讓你人在外地也能幫爸媽修手機
  \item 微信開"關懷模式"+"語音輸入"是長輩最實用的兩個功能
  \item SOS緊急求助和跌倒檢測，關鍵時刻救命的配置值得花五分鐘設好
\end{itemize}

\end{tcbsummary}
